\ifextended

\templateentry
{Diâmetro de uma árvore}
{}
{\texttt{treeDiam(n,g)} retorna a sequência de vértices no diâmetro da árvore. $O(n)$.}
{tree/tree-diameter.cpp}

\fi

\ifextended
\templateentry
{Matching em árvores}
{\texttt{treeMatching(g,root)} retorna o parceiro de cada vértice em $O(n)$.}
{Parceiro de cada vértice em um matching máximo; \texttt{-1} se livre. $O(n)$.}
{tree/tree-matching.cpp}

\else
\notebookentry{Matching em árvores}
{Processe os vértices de baixo para cima; pareie cada vértice não raiz com seu pai se ambos estiverem livres. Produz um matching máximo em $O(n)$.}
{}
\fi

\ifextended

\templateentry
{Isomorfismo de subárvores}
{}
{\texttt{TreeIso t}; \texttt{t.classify(g,r=0)}: classes exatas de isomorfismo de subárvores não ordenadas e sem rótulos, sob a raiz $r$; árvore não vazia. Reutilize a mesma instância para comparar árvores diferentes. Para $n$ vértices e $K$ classes já registradas: $O(n\log(n+K))$ tempo, memória proporcional às assinaturas armazenadas mais $O(n)$.}
{tree/isomorphism.cpp}

\fi

\templateentry
{LCA}
{\texttt{LCA lca(g,root)} e \texttt{lca.lca(a,b)}; construção $O(n \log n)$, consulta $O(1)$.}
{\texttt{LCA lca(g,root)}, \texttt{lca.lca(a,b)}. Construção $O(n\log n)$; consulta $O(1)$.}
{tree/lca.cpp}

\templateentry
{HLD}
{\texttt{HLD(g,r)}: árvore não vazia. Subárvore em $[\mathtt{tin}[u],\mathtt{tout}[u])$; \texttt{path(u,v)} retorna pares de extremos \textbf{inclusivos}, na ordem $u\to v$ (podem ser decrescentes). Construção $O(n)$; $O(\log n)$ intervalos por caminho.}
{\texttt{HLD(g,r)}: árvore não vazia. Subárvore em $[\mathtt{tin}[u],\mathtt{tout}[u])$; \texttt{path(u,v)} retorna pares de extremos \textbf{inclusivos}, na ordem $u\to v$ (podem ser decrescentes). Construção $O(n)$; $O(\log n)$ intervalos por caminho.}
{tree/hld.cpp}

\templateentry
{Árvore virtual}
{\texttt{virtree(...,v)}: arestas \texttt{\{pai,filho\}}. \texttt{v} vira o fechamento por LCA, ordenado por \texttt{tin} e sem duplicatas; \texttt{tin,tout} da mesma raiz, subárvore $[\mathtt{tin},\mathtt{tout})$. Para $k$ vértices e LCA $O(1)$: $O(k\log k)$.}
{\texttt{virtree(...,v)}: arestas \texttt{\{pai,filho\}}. \texttt{v} vira o fechamento por LCA, ordenado por \texttt{tin} e sem duplicatas; \texttt{tin,tout} da mesma raiz, subárvore $[\mathtt{tin},\mathtt{tout})$. Para $k$ vértices e LCA $O(1)$: $O(k\log k)$.}
{tree/virtual-tree.cpp}

\templateentry
{Árvore de reconstrução de Kruskal}
{\texttt{KruskalTree(n,e)}, arestas \texttt{(peso,u,v)} não orientadas. Folhas $[0,n)$; campos: \texttt{ch} filhos, \texttt{w} peso, \texttt{cnt} número de folhas, \texttt{p} pai; raízes em \texttt{roots}, pai $-1$. Para folhas distintas conexas, peso do LCA = menor limiar que as conecta. Floresta desconexa permitida. $O(n+m\log m)$, até $2n-1$ nós.}
{\texttt{KruskalTree(n,e)}, arestas \texttt{(peso,u,v)} não orientadas. Folhas $[0,n)$; campos: \texttt{ch} filhos, \texttt{w} peso, \texttt{cnt} número de folhas, \texttt{p} pai; raízes em \texttt{roots}, pai $-1$. Para folhas distintas conexas, peso do LCA = menor limiar que as conecta. Floresta desconexa permitida. $O(n+m\log m)$, até $2n-1$ nós.}
{tree/kruskal-tree.cpp}

\ifextended

\templateentry
{Decomposição em árvore de centróides}
{\texttt{centroidDecomp(g)}: pais na árvore de centróides, raiz com $-1$. Requer árvore não vazia. $O(n\log n)$.}
{\texttt{centroidDecomp(g)}: pais na árvore de centróides, raiz com $-1$. Requer árvore não vazia. $O(n\log n)$.}
{tree/centroid-decomp.cpp}

\templateentry
{Histograma de distâncias em árvore}
{}
{\texttt{treeDistHist(g)}: \texttt{ans[d]} conta pares não ordenados de vértices distintos à distância \texttt{d}. Implemente \texttt{conv(vector<ll>,\allowbreak vector<ll>)} usando fft com \texttt{double}. $O(n \log^2 n)$.}
{tree/dist-hist.cpp}

\fi

\templateentry
{Link-cut tree}
{\texttt{LinkCut<S,op,e> t(n)}: vértices \textbf{1 a n}, sem arestas; monoide como na segtree. \texttt{set(u,x)}; \texttt{query(u,v)} agrega $u\to v$, extremos inclusos, requer conexão; aceita \texttt{op} não comutativa. \texttt{link/cut(u,v)}: falso se criaria ciclo/aresta ausente. $O(\log n)$ amortizado com \texttt{op} e cópias $O(1)$; memória $O(n)$.}
{\texttt{LinkCut<S,op,e> t(n)}: floresta sem arestas, vértices \textbf{1 a n}. Monoide como na segtree. \texttt{set(u,x)}; \texttt{query(u,v)} agrega o caminho na ordem $u\to v$, incluindo ambos os extremos, e requer conexão; aceita operação não comutativa. \texttt{link/cut(u,v)}: falso se criaria ciclo/aresta ausente. Operações $O(\log n)$ amortizado com \texttt{op} e cópias $O(1)$; memória $O(n)$. Agrega apenas caminhos, não subárvores.}
{tree/link-cut.cpp}
